% !TEX root = branch_time_ru.tex
\documentclass[12pt,a4paper]{article}
\usepackage{cmap}
\usepackage[margin=1in]{geometry}
\usepackage{amsmath,amssymb,amsfonts}
\usepackage[utf8]{inputenc}
\usepackage[T2A]{fontenc}
\usepackage[russian]{babel}
\usepackage[hidelinks]{hyperref}
\usepackage{comfortaa}
\renewcommand*\familydefault{\sfdefault}
\usepackage[table,x11names]{xcolor}
\usepackage{booktabs}
\usepackage{array}
\usepackage{graphicx}

\makeatletter
\renewcommand{\@makefnmark}{%
  \hbox{\@textsuperscript{\normalfont\color{white}\@thefnmark}}}
\renewcommand{\@makefntext}[1]{%
  \parindent=1em\color{white}\raggedright
  \hb@xt1.8em{\hss\@makefnmark\hskip0.3em}#1}
\renewcommand{\footnoterule}{%
  \kern-3\p@\color{white}\hrule width 0.4\columnwidth height 0.4\p@
  \kern2.6\p@}
\makeatother

\title{\textbf{Собственное время, зависящее от ветви}}
\author{Олег Понфиленок\thanks{Независимый исследователь; ponfil@gmail.com}}
\date{}

\begin{document}

\pagecolor{black}
\color{white}
\makeatletter
\def\normalcolor{\color{white}}
\makeatother

\maketitle

\begin{abstract}
Мы предлагаем привязывать фундаментальный масштаб гравитационно индуцированной «размытости» пространства-времени к \emph{собственному времени}, а не к пространственной длине, и постулируем, что при ансамблевом (статистическом) прочтении квантового состояния, в котором каждая ветвь несёт свою мировую линию, флуктуация собственного времени \emph{статистически независима по ветвям}. Две посылки вместе влекут, что в относительную фазу входит \emph{полная} энергия покоя ветви $mc^2$, а не только разность энергий между ветвями. Это даёт функционал декогеренции $V/V_0=\exp(-\chi)$ с $\chi=(mc^2/\hbar)^2(\xi\,t_P^2\,t)^{2/3}$, где $t$~--- время пребывания в суперпозиции, а $\xi=\mathcal O(1)$~--- числовой коэффициент; форма скейлинга $\chi\propto m^2 t^{2/3}$ следует из постулатов, тогда как общий коэффициент несёт $\mathcal O(1)$-неопределённость. Скейлинг $\chi\propto m^2 t^{2/3}$ совпадает с квадратичным по массе поведением, ожидаемым для конъектурных флуктуаций пространства-времени, но механизм резко отличается от коллапса Диоши--Пенроуза и от стандартной модели Каройхази. При сопоставлении с интерферометрией наночастиц натрия венской группы (2026; $m\approx1{,}7\times10^5$~Да, $t\sim 6$--$12$~мс) данные накладывают условия на $\xi\lesssim0{,}3$. Решающий тест близок: при $m\sim 1$~МДа и сниженной скорости $\chi\gg1$ и интерференция исчезает полностью, тогда как стандартная квантовая механика предсказывает полный контраст.
\end{abstract}

\noindent\textbf{Ключевые слова:} собственное время; независимость по ветвям; неопределённость Каройхази; декогеренция; интерферометрия материи; теорема Марголуса--Левитина

%%% ===================================================================
\section{Введение: ансамблевая онтология и время}
%%% ===================================================================

Статус волновой функции~--- онтический объект или вычислительный учёт ансамбля возможностей~--- остаётся нерешённым. Мы принимаем второе, ансамблевое прочтение~\cite{ballentine1970,home1992}: $\lvert\psi\rangle=\sum_i c_i\lvert i\rangle$ каталогизирует различные \emph{варианты} системы, а квадраты амплитуд~--- относительные веса этих вариантов. Это намеренно минимальная онтология, но она несёт немедленное следствие, как только в игру вводится гравитация.

В реляционном взгляде на время часы~--- это физическая подсистема, а «время»~--- корреляция внутри полного состояния~\cite{pagewootters}. Если каждая ветвь есть подлинный вариант системы, то она несёт свою мировую линию и, следовательно, своё собственное время. Пенроуз приходит к структурно тождественному утверждению с противоположной стороны: каждая компонента суперпозиции источает своё пространство-время, и «поскольку нет канонического способа связать временные параметры в двух пространствах-временах, возникает фундаментальная неопределённость оператора сдвига во времени»~\cite{penrose1996}. Согласованно с этим, собственное время накапливается вдоль мировой линии и не задаётся одной общей шкалой для различных четырёхмерных путей: в интерферометре с разными траекториями нельзя без дополнительного постулата навесить на все ветви один общий $\tau$. Пока система в суперпозиции путей, интерференция требует когерентной относительной фазы между ними; статистически независимые по ветвям флуктуации $\delta t$ (постулат~П2 ниже) вводят случайный сдвиг этой фазы и ограничивают видность сверху, $V/V_0\le e^{-\chi}$ (см.~\eqref{eq:chi})~--- это предсказание предела интерференции при собственном времени, зависящем от пути, а не отрицание амплитудной суперпозиции. После регистрации which-path фиксируется один сценарий и один $\tau$ вдоль реализованной линии; $|c_i|^2$ остаются весами исходов. Механизм ниже~--- чистая дефазировка между ветвями: популяции не перераспределяются, which-path-информация не записывается.

Настоящее предложение отличается от пенроузовского в двух отношениях. Во-первых, мы берём флуктуирующим объектом само собственное время; пространственная неопределённость длины Каройхази тогда~--- производное следствие, а не первооснова. Операционально это естественно: с 1983 года метр \emph{определяется} через секунду, так что неопределённость длины вторична по отношению к неопределённости времени. Во-вторых~--- и здесь модель становится предсказательной~--- мы постулируем, что флуктуации собственного времени разных ветвей статистически \emph{независимы}. Один этот шаг воскрешает полную энергию покоя каждой ветви как физически активную фазу, порождая закон декогеренции, квадратичный по массе и уже находящийся на границе лабораторной достижимости.

%%% ===================================================================
\section{Два постулата}
%%% ===================================================================

Введём ключевые обозначения:
\begin{itemize}
\item $t$~--- время пребывания системы в суперпозиции (интервал собственного времени вдоль мировой линии);
\item $\delta t$~--- масштаб неопределённости собственного времени за интервал $t$ (характерная величина случайного сдвига $\delta\tau$ вдоль мировой линии);
\item $t_P=\sqrt{\hbar G/c^5}$~--- планковское время;
\item $\ell_P=\sqrt{\hbar G/c^3}$~--- планковская длина;
\item $m$~--- масса частицы (или кластера) в каждой ветви;
\item $\omega_C=mc^2/\hbar$~--- комптоновская частота;
\item $\xi=\mathcal O(1)$~--- безразмерный числовой коэффициент.
\end{itemize}

\paragraph{(П0) Структура ветвления.}
Ветви~--- это компоненты состояния, различающиеся \emph{мировой линией носителя}: неперекрывающиеся волновые пакеты центра масс (which-path); для чисто внутренних степеней свободы ветвление задаётся собственным базисом энергии. Один связный волновой пакет несёт одну мировую линию и одни часы; ветвление наступает, когда эволюция расщепляет состояние на различимые (взаимно ортогональные по носителю) компоненты, и прекращается при их слиянии или регистрации which-path. Подчеркнём: это дополнительный \emph{структурный} постулат, а не следствие ансамблевой онтологии разд.~1. Дефазировка вида $\rho_{ab}\to e^{-\chi}\rho_{ab}$ определена лишь относительно выделенного разложения на ветви~--- одно и то же состояние в разных базисах содержит разное число компонент,~--- и без~(П0) закон~\eqref{eq:chi} не задаёт отображения состояний.

\paragraph{(П1) Первичность и величина флуктуации собственного времени.}
Накопленная неопределённость интервала собственного времени $t$, измеренного вдоль мировой линии, равна
\begin{equation}
\boxed{\delta t^3=\xi\,t_P^2\,t,}
\label{eq:dt}
\end{equation}
с $\xi=\mathcal O(1)$. Уравнение~(\ref{eq:dt})~--- временная форма соотношения Каройхази $\Delta s^3=\ell_P^2 s$~\cite{karolyhazy1966}. Оно получается \emph{внутренне из часов}, без обращения к стержням, аргументом Залекера--Вигнера~\cite{saleckerwigner} / Нга--ван Дама~\cite{ngvandam1995,ngvandam1994}: часы массы $m$, измеряющие $t$ с разрешением $\delta t$, подчиняются (квантовое расплывание) $\delta t^2\gtrsim\hbar t/(mc^2)$, откуда $m\gtrsim\hbar t/(c^2\delta t^2)$; одновременно требование, чтобы часы не были чёрной дырой, $\ell\gtrsim r_s=2Gm/c^2$, даёт $m\lesssim c^3\delta t/(2G)$. Исключая $m$, получаем $\delta t^3\gtrsim 2\,t_P^2\,t$, то есть уравнение~(\ref{eq:dt}) с $\xi=2$.

Подчеркнём, что гравитация существенна: без ограничения «не чёрная дыра» поле исчезает ($\delta t\to0$ при $m\to\infty$). Не менее существенна и собственная масса покоя: неопределённость $\delta t$ определена вдоль мировой линии с ненулевым собственным временем; без массы покоя собственного времени нет. У фотона $m=0$, $\mathrm{d}\tau=0$, и весь механизм к светоподобному носителю не применим. Показатель степени $1/3$ и планковский масштаб $t_P^2$ гравитационны по происхождению, хотя вывод нигде не привлекает пространственную длину; тот же кубический закон есть голографический масштаб неопределённости пространства-времени~\cite{ng2005foam}. Коэффициент $\xi$ на этом уровне остаётся неопределённым~--- это $\mathcal O(1)$-неопределённость.

\paragraph{(П2) Независимость по ветвям.}
Для двух ветвей $a,b$ (в смысле~П0) реализованные сдвиги собственного времени $\delta\tau_a,\delta\tau_b$~--- независимые случайные величины, каждая с дисперсией $\delta t^2$ из уравнения~(\ref{eq:dt}). Это содержательный постулат, мотивированный ансамблевой онтологией разд.~1: разные варианты не разделяют общих часов. Не менее важна обратная сторона постулата: \emph{внутри} одной ветви все составляющие разделяют одни часы~--- часы привязаны к ветви как целому, к мировой линии её центра масс, а не к отдельным частицам.

\paragraph{Одни часы на ветвь: центр масс как носитель энергии покоя.}
Выбор «одни часы на вариант» не техническая деталь~--- от него зависит сам закон $m^2$. В разложении гамильтониана связанной системы
\begin{equation}
H = Mc^2 + \frac{P^2}{2M} + H_{\mathrm{int}}
\label{eq:Hcm}
\end{equation}
полная энергия покоя $Mc^2$ сидит на \emph{одной} коллективной координате~--- центре масс; именно его мировая линия ветвится в интерферометре, и именно к ней привязаны часы ветви. Тогда для кластера из $N$ атомов в суперпозиции путей относительная фаза есть $\Delta\varphi=(Mc^2/\hbar)(\delta\tau_a-\delta\tau_b)$ и $\chi\propto M^2=(N m_{\mathrm{at}})^2$. Если бы, напротив, каждый атом нёс собственные независимые часы, дисперсии фаз складывались бы аддитивно и $\chi\propto N\,m_{\mathrm{at}}^2$~--- в $N\approx5000$--$10^4$ раз меньше для венских кластеров разд.~4 ($\chi\sim10^{-4}$ вместо $\sim0{,}4$), и предсказательное содержание модели исчезло бы. Квадратичный по массе закон, таким образом, эквивалентен утверждению, что носителем флуктуирующего собственного времени является ветвь как целое, а не её составляющие по отдельности.

\paragraph{Природа неопределённости: информационная, а не средовая.}
Существенно, что $\delta t$ здесь~--- не флуктуация внешнего метрического поля и не средовой шум, общий для близких траекторий. Это \emph{информационная} неопределённость самого собственного времени: при голографическом пределе на плотность информации ни одна величина, включая показание любых часов, не определена с бесконечной точностью, а кубический закон~\eqref{eq:dt} есть как раз голографический масштаб такой неопределённости~\cite{ng2005foam}. Различие принципиально для~(П2). Средовой источник~--- например, разность гравитационных потенциалов на плечах интерферометра~--- дал бы \emph{скоррелированный} сдвиг, тем более близкий к общему, чем ближе ветви, и в относительной фазе он сократился бы; но это систематический сдвиг, а не шум. Голографическая же неопределённость не наводится ничем, что ветви разделяют: по~(П0) она привязана к ветви как целому~--- к мировой линии её центра масс,~--- а не к точкам пространства, поэтому пространственная близость ветвей сама по себе не создаёт корреляции их часов. Привязка именно к дискретным ветвям, а не к положениям, существенна: шум, независимый для сколь угодно близких \emph{положений}, означал бы случайное поле фазы с неограниченным градиентом~--- то есть силу и диффузию импульса; шум, независимый для дискретных \emph{ветвей}, есть чистая дефазировка без силы (разд.~4.1). В ансамблевом прочтении, где каждая ветвь есть самостоятельный вариант со своей мировой линией, величины $\delta\tau_a,\delta\tau_b$ оттого статистически независимы. Тем самым~(П2)~--- не дополнительное предположение о среде, а следствие информационной первичности времени в применении к различным вариантам. Несводимым остаётся лишь сама посылка, что голографический бюджет информации выделяется \emph{поветвенно}, а не как единый общий регистр; именно она несёт предсказательную нагрузку модели.

%%% ===================================================================
\section{Функционал декогеренции}
%%% ===================================================================

Фаза покоя ветви набегает на комптоновской частоте $\omega_C=mc^2/\hbar$. Согласно (П2), относительная фаза двух ветвей равной массы есть $\Delta\varphi=\omega_C(\delta\tau_a-\delta\tau_b)$~--- глобальная фаза покоя, которая сократилась бы при \emph{общих} часах, но выживает здесь, потому что часы независимы. Её дисперсия
\begin{equation}
\mathrm{Var}(\Delta\varphi)=2\,\omega_C^2\,\delta t^2 .
\end{equation}
Усреднение $\langle e^{i\Delta\varphi}\rangle$ по гауссовой фазе снижает видность интерференции на множитель $\exp(-\mathrm{Var}/2)$; множитель $2$ от двух ветвей и множитель $1/2$ от гауссова усреднения сокращаются
\begin{equation}
\boxed{\;\frac{V}{V_0}=e^{-\chi},
\qquad
\chi=\omega_C^2\,\delta t^2
=\left(\frac{mc^2}{\hbar}\right)^{2}\!\bigl(\xi\,t_P^2\,t\bigr)^{2/3}.\;}
\label{eq:chi}
\end{equation}
Форма скейлинга $\chi\propto m^2 t^{2/3}$ следует из постулатов и составляет центральное проверяемое предсказание; общий же коэффициент несёт $\mathcal O(1)$-неопределённость (через $\xi^{2/3}$, трактовку $\delta t$ и выбор эффективного $t$; подробнее в разд.~4). Квадратичная по массе зависимость соответствует общему ожиданию для конъектурных флуктуаций пространства-времени в литературе по интерферометрии материи~\cite{kialka2022}.

\paragraph{Планковская граница масс.}
Приравнивая одночастичное время декогеренции $t_{\mathrm{dec}}=t_C^3/t_P^2$ (время $\chi=1$ для одного носителя) собственному комптоновскому времени носителя $t_C=\hbar/mc^2$, получаем $t_C=t_P$, то есть
\begin{equation}
m=m_P=\sqrt{\hbar c/G}\approx 2{,}2\times10^{-8}\;\text{кг}.
\end{equation}
Модель воспроизводит стандартное ожидание, что квантово-классическая граница лежит на планковской массе. Иначе: масса (энергия покоя $mc^2$), расщеплённая между когерентными ветвями, ограничена сверху $m_P$~--- при $m>m_P$ время декогеренции $t_{\mathrm{dec}}\ll t_C$ и когерентность между ветвями теряется на временах, несоизмеримо меньших любого макроскопического процесса. Кот Шрёдингера ($M\sim1$~кг $\gg m_P$) есть предельный пример: суперпозиция «жив/мёртв» амплитудно записываема, но интерференция между ветвями гаснет мгновенно; в ансамблевом прочтении остаётся статистическая смесь с весами $|c_i|^2$~--- чистая дефазировка, без перераспределения популяций (в отличие от объективного коллапса) и без спонтанного излучения, которое предсказывают модели Диоши--Пенроуза~\cite{donadi2021}.

\paragraph{Что входит в массу ветви $m$.}
Уточним, какая масса входит в $m$ закона~(\ref{eq:chi}). Критерий касается \emph{когерентных ветвей}, а не пространства: вклад даёт энергия покоя \emph{только тех степеней свободы, чьё физическое состояние различается между интерферирующими ветвями} (которые запутаны с меткой ветви). Лишь они расщеплены по ветвям и обретают независимые часы (П2). Масса же, входящая \emph{общим тензорным множителем}~--- в одном и том же состоянии во всех ветвях,~--- не расщеплена: у одного физического объекта одни часы, его сдвиг $\delta\tau$ в ветвях один и тот же, и его фаза покоя сокращается в относительной фазе. Формально для $|\psi\rangle=|C\rangle\otimes(\alpha|0\rangle+\beta|1\rangle)$ общий множитель $|C\rangle$ (несущая среда, прибор, покоящееся в одном состоянии ядро) даёт глобальную фазу и выпадает; ветвится лишь $(\alpha|0\rangle+\beta|1\rangle)$. Поэтому «масса присутствует в системе» и даже «система в суперпозиции» $\neq$ «масса различается между ветвями». Полная масса покоя входит в $m$ \emph{только} тогда, когда по ветвям различается состояние самого носителя~--- ярчайший случай суперпозиция положений центра масс (which-path), как в интерферометрии материи (разд.~4). Внутреннее же состояние (спин, гиперточный или оптический переход) оставляет носитель общим множителем, и его масса покоя в $m$ не входит~--- входит лишь энергия реально различающейся степени свободы. Эмпирически это подтверждается сверхпроводящими резонаторами, хранящими одиночный фотон $\sim\!34$~мс при $\gtrsim10^{23}$ электронах в стенках~\cite{milul2023cavity}: будь масса конденсата расщеплена по ветвям, когерентность гасла бы за наносекунды.

%%% ===================================================================
\section{Сопоставление с экспериментом}
%%% ===================================================================

\subsection{Интерферометрия материи: единственный зонд}

По построению~(П0)/(П2) шум привязан к дискретным ветвям, а не к точкам пространства, поэтому механизм есть чистая дефазировка фазы покоя: нет позиционно-зависимого стохастического потенциала, значит нет силы, нет диффузии импульса и нет спонтанного излучения. Подчеркнём, что это свойство обеспечено именно постулатом ветвления~(П0), а не автоматично: шум, привязанный к положениям, наводил бы силу. Как следствие, модель \emph{прозрачна} для неинтерферометрических ограничений, держащих Диоши--Пенроуз: подземный поиск спонтанного излучения~\cite{donadi2021}, предел Majorana Demonstrator~\cite{majorana2022} и рентгеновские тесты~\cite{piscicchia2024} к ней неприменимы. Механизм не наводит и детектируемой нестабильности хода атомных часов: эффект существует лишь \emph{между} ветвями суперпозиции, тогда как одиночные (несуперпонированные) часы фазы покоя не накапливают разностно, так что их ход остаётся незатронутым. Поэтому интерферометрия~--- единственный лабораторный зонд.

\subsection{Наночастицы натрия: венская группа (2026)}

\paragraph{Параметры установки.}
Сопоставляем уравнение~(\ref{eq:chi}) с интерферометром наночастиц натрия венской группы~\cite{arndt2025,arndt_arxiv}:
\begin{itemize}
\item кластеры из $5000$--$10000$ атомов Na с центром распределения масс $m\approx 172$~кДа ($m\approx 2{,}86\times10^{-22}$~кг);
\item скорость $v\approx 160$~м/с;
\item три решётки периода $d=133$~нм на расстоянии $L=0{,}983$~м;
\item измеренная видность полос $V=0{,}10\pm0{,}01$ (второй прогон $0{,}08\pm0{,}01$) при $172$~кДа, описываемая квантовой моделью при едином глобальном множителе обычных потерь контраста $0{,}78$ (юстировка, гравитационное и вращательное усреднение фазы, вибрации, тепловая и столкновительная декогеренция);
\item рекордная макроскопичность $\mu=15{,}45$~\cite{macroscopicity,arndt2025}.
\end{itemize}
Релевантное время суперпозиции~--- тальботовский пролёт $t=L/v\approx 6{,}1$~мс (одна секция) вплоть до $2L/v\approx 12{,}3$~мс.

\paragraph{Численная оценка и ограничение на $\xi$.}
С $\omega_C=mc^2/\hbar=2{,}43\times10^{29}$~с$^{-1}$ (при $m=172$~кДа) и $t_P^2=2{,}91\times10^{-87}$~с$^2$ уравнение~(\ref{eq:chi}) даёт
\begin{equation}
\chi\approx 0{,}41\,\xi^{2/3}\quad (t=L/v\approx 6{,}1~\text{мс}),
\qquad
\chi\approx 0{,}65\,\xi^{2/3}\quad (t=2L/v\approx 12{,}3~\text{мс}).
\end{equation}
При выводном $\xi\approx2$ (П1) это $\chi\approx 0{,}64$--$1{,}02$, то есть дефицит видности $48$--$64\%$. Поскольку наблюдаемые полосы описываются квантовой моделью с множителем обычных потерь $0{,}78$, дополнительный множитель $e^{-\chi}$ совместим с данными лишь при $e^{-\chi}\gtrsim0{,}78$, то есть $\chi\lesssim0{,}25$, откуда
\begin{equation}
\boxed{\;\xi\lesssim 0{,}3.\;}
\end{equation}
Подчеркнём, однако, что сам коэффициент в $\chi$ эвристически не закреплён: $\xi$ (вывод даёт $\approx2$, но честно «$\mathcal O(1)$»), трактовка $\delta t$ (СКО / полуширина / минимальная неопределённость), выбор эффективного $t$ и перевод $\delta t$ в фазовый критерий (порядка одного радиана против $\pi$) совокупно дают неопределённость множителя в $\chi$ вплоть до $\sim\pi^2$. Сам закон $V/V_0=e^{-\chi}$ при этом точен (гауссова дефазировка). Поэтому датум при одной массе модель не \emph{исключает}, а лишь \emph{ограничивает} эффективный коэффициент; формально вклад $e^{-\chi}$ при фиксированной массе вырожден с глобальным множителем $0{,}78$ и от него не отделяется. На текущих данных модель ни безопасно согласуется, ни чисто отвергнута.

\begin{table}[ht]
\centering
\arrayrulecolor{white}
\caption{Дефицит видности $1-V/V_0$ для интерферометра натрия при разных $\xi$.}
\label{tab:visibility}
\begin{tabular}{l c c c c}
\hline
 & & \multicolumn{3}{c}{дефицит видности $1-V/V_0$}\\
Режим & $t$ & $\xi=0{,}3$ (допуск.) & $\xi=1$ & $\xi=2$ (вывод) \\
\hline
1 секция & $6{,}1$~мс & $\approx 17\%$ & $\approx 33\%$ & $\approx 47\%$ \\
2 секции & $12{,}3$~мс & $\approx 25\%$ & $\approx 47\%$ & $\approx 64\%$ \\
\hline
\end{tabular}
\end{table}
Эксперимент поглощает в глобальный множитель $0{,}78$ ($V_{\rm obs}=V_{\rm QM}\times0{,}78$) потерю контраста до $\approx\!22\%$, поэтому дефициты $\lesssim\!22\%$ им допускаются. Допускаемое данными $\xi\approx0{,}3$ даёт $17$--$25\%$~--- на самой границе этого бюджета, то есть эксперимент такую модель \emph{разрешает}; выводное $\xi\approx2$ ($47$--$64\%$) исключается. Таким образом, текущие данные не отвергают модель, но прижимают её коэффициент к $\xi\lesssim0{,}3$.

\subsection{Решающий ближайший тест}

Уже измеренная при $0{,}4$--$1$~МДа высокая видность $V=0{,}66\pm0{,}09$~\cite{arndt2025} механизм \emph{не} тестирует: она снята в режиме коротких длин волн $\lambda_{dB}\lesssim3$~фм, $L\le L_T$, где квантовое и классическое предсказания совпадают (геометрическая оптика), то есть подлинной квантовой когерентности, которую дефазирует наш механизм, там нет. Различающий тест требует квантового режима при большой массе. Та же программа нацелена на $m\sim 1$~МДа с частицами, замедленными до $v\approx 25$~м/с~\cite{arndt2025}. Тогда $t\approx 2L/v\approx 79$~мс, $\omega_C\approx 1{,}4\times10^{30}$~с$^{-1}$, и уравнение~(\ref{eq:chi}) даёт
\begin{equation}
\chi\approx 50\text{--}75,
\qquad
V/V_0=e^{-\chi}\sim10^{-20},
\end{equation}
то есть \emph{полное} гашение интерференции, устойчивое к любой неопределённости коэффициента (даже деление $\chi$ на $\sim\pi^2$ оставляет $\chi\sim5$--$7$). Стандартная квантовая механика предсказывает полосы полного контраста на этой массе. Исходы, следовательно, бинарны и близки во времени: наблюдение полос при $\sim\!1$~МДа сразу фальсифицирует модель; их исчезновение со скейлингом $V\propto\exp[-(mc^2/\hbar)^2(\xi t_P^2 t)^{2/3}]$ было бы свидетельством в её пользу.

\subsection{Отличающая сигнатура}

Один лишь дефицит видности не диагностичен, так как тепловое излучение, столкновения с остаточным газом и разброс скоростей тоже снижают контраст. Сигнатура настоящего механизма~--- его скейлинг $\chi\propto m^2 t^{2/3}$ вместе с \emph{невосприимчивостью к экранировке}: его нельзя уменьшить улучшением вакуума или охлаждением частицы. Экспериментальный различитель~--- сохраняется ли остаточный дефицит после подавления всех известных средовых каналов и следует ли он закону $m^2$.

%%% ===================================================================
\section{Предел Марголуса--Левитина: максимальное число ортогональных
операций в окне когерентности}
%%% ===================================================================

Время декогеренции модели~--- момент, когда $\chi=1$ и видность падает до $e^{-1}$. Из уравнения~(\ref{eq:chi})
\begin{equation}
t_{\mathrm{dec}}=\frac{t_C^3}{\xi\,t_P^2},
\qquad
t_C=\frac{\hbar}{mc^2}=\omega_C^{-1}.
\label{eq:tdec}
\end{equation}
Это максимальное время, в течение которого ветви ещё сохраняют относительную фазу порядка единицы; дальше интерференция гаснет.

Теорема Марголуса--Левитина~\cite{margolus1998} ограничивает скорость ортогональной эволюции средней энергией системы над основным уровнем $E$: минимальное время перехода в ортогональное состояние
\begin{equation}
\tau_\perp\ge\frac{\pi\hbar}{2E},
\end{equation}
откуда максимальная частота ортогональных операций
\begin{equation}
\nu_{\max}=\frac{2E}{\pi\hbar}.
\label{eq:nuML}
\end{equation}
Для носителя с энергией покоя $E=mc^2$ это $\nu_{\max}=2\omega_C/\pi$. За окно когерентности~(\ref{eq:tdec}) максимальное число ортогональных тактов~---
\begin{equation}
N_{\max}=\nu_{\max}\,t_{\mathrm{dec}}
=\frac{2\omega_C}{\pi}\cdot\frac{t_C^3}{\xi\,t_P^2}
=\frac{2}{\pi\xi}\left(\frac{t_C}{t_P}\right)^{2}.
\end{equation}
Поскольку $t_C/t_P=m_P/m$, получаем
\begin{equation}
\boxed{\;N_{\max}=\frac{2}{\pi\xi}\left(\frac{m_P}{m}\right)^{2}.\;}
\label{eq:Nmax}
\end{equation}
Максимальное число ортогональных операций в окне когерентности зависит только от отношения массы к планковской и от $\xi$; комптоновская частота сокращается. При $m=m_P$ уравнение~(\ref{eq:Nmax}) даёт $N_{\max}=2/(\pi\xi)$; ровно одна ортогональная операция получается при
\begin{equation}
\xi=\frac{2}{\pi}\approx0{,}637,
\end{equation}
то есть на планковской массе окно когерентности вмещает ровно один ортогональный такт~--- та же граница, что уже воспроизведена в разд.~3 из $t_{\mathrm{dec}}=t_C$ (при $\xi=\mathcal O(1)$).

\paragraph{Численные оценки.}
Таблица~\ref{tab:Nmax} приводит $N_{\max}$ при допускаемом данными $\xi\lesssim0{,}3$, при $\xi=1$ и при $\xi=2/\pi$ (ровно $N_{\max}=1$ на $m_P$):
\begin{table}[ht]
\centering
\arrayrulecolor{white}
\caption{Максимальное число ортогональных операций $N_{\max}$ за окно когерентности $\chi=1$.}
\label{tab:Nmax}
\begin{tabular}{l c c c c}
\hline
Носитель & $m$ & $N_{\max}$ ($\xi=0{,}3$) & $N_{\max}$ ($\xi=2/\pi$) & $N_{\max}$ ($\xi=1$) \\
\hline
электрон & $9{,}1\times10^{-31}$~кг & $\sim10^{45}$ & $\sim10^{44}$ & $\sim10^{44}$ \\
протон & $1{,}7\times10^{-27}$~кг & $\sim10^{38}$ & $\sim10^{38}$ & $\sim10^{38}$ \\
Na, $172$~кДа & $2{,}86\times10^{-22}$~кг & $\sim10^{28}$ & $\sim10^{27}$ & $\sim10^{27}$ \\
$1$~МДа & $1{,}66\times10^{-21}$~кг & $\sim10^{26}$ & $\sim10^{26}$ & $\sim10^{26}$ \\
$m_P$ & $2{,}2\times10^{-8}$~кг & $\sim2$ & $1$ & $2/\pi\approx0{,}64$ \\
\hline
\end{tabular}
\end{table}
Для лабораторных масс окно когерентности ещё огромно по меркам Марголуса--Левитина: даже мегадальтонный кластер формально допускает максимум $\sim10^{26}$ ортогональных тактов до $\chi=1$. Смысл оценки~--- не в том, что такой ресурс достижим на практике, а в том, что модель накладывает \emph{конечный верхний} бюджет ортогональной эволюции на каждую массу: $N_{\max}\propto m^{-2}$. Чем тяжелее расщеплённый носитель, тем меньше максимальное число ортогональных операций, умещающихся, пока ветви ещё когерентны. На планковской массе этот максимум схлопывается до $\mathcal O(1)$~--- согласованно с тем, что там $t_{\mathrm{dec}}\sim t_C\sim t_P$.

\paragraph{Максимум операций за заданное время: голографический скейлинг.}
Можно поставить вопрос иначе: сколько ортогональных операций вообще возможно за заданное время $t$, если массу носителя разрешено выбирать? Скорость Марголуса--Левитина $N=(2/\pi)\,t/t_C$ растёт с массой, но условие сохранения когерентности $t_{\mathrm{dec}}\ge t$ ограничивает массу сверху: $t_C\ge(\xi\,t_P^2\,t)^{1/3}$. Максимум достигается на границе, где $t_C=(\xi\,t_P^2\,t)^{1/3}$, и масса из ответа выпадает:
\begin{equation}
\boxed{\;N_{\max}(t)=\frac{2}{\pi\,\xi^{1/3}}\left(\frac{t}{t_P}\right)^{2/3}.\;}
\label{eq:Nmax_t}
\end{equation}
За одну секунду это $\approx4{,}5\times10^{28}$ операций (при $\xi=1$; оптимальная масса $\sim5\times10^{4}$~а.е.м.). Показатель $2/3$ здесь не случаен: тот же скейлинг $N\propto(t/t_P)^{2/3}$ Нг получает как голографический предел на вычислительную мощность «компьютера из пространственно-временной пены» и чёрнодырного квантового компьютера~\cite{ng2005foam}~--- из того же кубического закона $\delta t^3\sim t_P^2\,t$, что лежит в основе постулата~(П1). Модель, таким образом, автоматически воспроизводит голографическую границу вычислений: когерентная ортогональная эволюция в ней не может обгонять бюджет, задаваемый пространственно-временной пеной.

%%% ===================================================================
\section{Пределы квантовых вычислений: шум между ветвями логического регистра}
%%% ===================================================================

Постулаты (П0)--(П2) действуют на любую суперпозицию, включая вычислительную суперпозицию квантового регистра: по~(П0) ветвление для внутренних степеней свободы задаётся собственным базисом энергии, для регистра~--- вычислительным базисом кубитов. Применим критерий массы ветви (разд.~3). Ветви алгоритма $|0\rangle,|1\rangle$ различаются \emph{состоянием кубита}, а не положением массы покоя носителя: масса чипа, атомов и резонаторов входит общим тензорным множителем и из относительной фазы выпадает. Физически активной остаётся лишь энергия реально различающейся степени свободы~--- энергия перехода кубита $\hbar\omega_q$. Тогда случайная относительная фаза между парой ветвей за время счёта $t$ имеет стандартное отклонение
\begin{equation}
\sigma(t)=\sqrt{2}\,\omega_q\,\delta t
=\sqrt{2}\,\omega_q\bigl(\xi\,t_P^2\,t\bigr)^{1/3}.
\label{eq:sigmaq}
\end{equation}
Для сверхпроводящего кубита ($\omega_q\approx2\pi\times5$~ГГц) и $t=1$~с уравнение~\eqref{eq:sigmaq} даёт $\delta t\approx1{,}4\times10^{-29}$~с и $\sigma\approx6\times10^{-19}$~рад: механизм действует и здесь, но как чрезвычайно слабая дефазировка. Существенно, что дисперсия~\eqref{eq:sigmaq} \emph{попарная} и от числа ветвей не зависит: в интерференцию входят лишь разности независимых $\delta\tau$, а экстремальный разброс фаз по $2^n$ ветвям растёт лишь как $\sigma\sqrt{2\ln 2^{\,n}}\propto\sigma\sqrt{n}$. Для пары ветвей, различающихся состоянием $d$ кубитов, $\chi\propto d^{\,2}$; в худшем случае, когда различаются все $n$ кубитов, формальный предел $n\,\omega_q\,\delta t\lesssim1$ даёт $n\lesssim10^{18}$~--- далеко за горизонтом практики.

\paragraph{Пол различимости фаз: точный и приближённый QFT.}
Величина $\sigma$ задаёт \emph{пол} для алгоритмических фазовых поворотов: гейт на угол $\theta\lesssim\sigma$ неотличим от собственного шума модели и не может нести информацию. Точное квантовое преобразование Фурье на $n$ кубитах содержит контролируемые повороты вплоть до $2\pi/2^{\,n}$, и уже при $n\gtrsim64$ эти углы опускаются ниже пола $\sim10^{-19}$~рад: точный QFT на больших регистрах модель запрещает. Однако точный QFT в реализациях и не используется: по Копперсмиту~\cite{coppersmith1994} все повороты с $k>\mathcal O(\log(n/\varepsilon))$ можно \emph{исключить из схемы} целиком, потеряв лишь $\varepsilon$ вероятности успеха; минимальный оставшийся угол $\sim2\pi\varepsilon/n$ полиномиально, а не экспоненциально мал. Для алгоритма Шора~\cite{shor1997} при $n=10^4$, $\varepsilon=10^{-2}$ это $\sim10^{-5}$~рад~--- на тринадцать порядков выше пола~\eqref{eq:sigmaq}. Фазово-оценочные алгоритмы, считывающие ответ положением интерференционного пика, механизм поэтому не ограничивает: он запрещает ровно те экспоненциально мелкие повороты, которые из схем и так исключаются.

\paragraph{Амплитудное усиление: алгоритм Гровера.}
Контрастен случай амплитудного усиления~\cite{grover1996}. Здесь QFT нет и усекать нечего: каждая итерация поворачивает вектор состояния на фиксированный угол $2\theta\approx2/\sqrt{N}$, заданный геометрией задачи, и все $\sim(\pi/4)\sqrt{N}$ итераций~--- одно непрерывное когерентное накопление без промежуточных считываний. Мелкий поворот за шаг неустраним: он и есть алгоритм. Для криптографических размеров ($N=2^{256}$) сигнальный шаг $2^{-127}\approx6\times10^{-39}$~рад лежит на двадцать порядков ниже пола~\eqref{eq:sigmaq}, а накопленная за $2^{128}$ итераций дефазировка астрономична при любом ненулевом шуме~--- хотя практически раньше убивает само время счёта. Это согласуется с известной структурной хрупкостью алгоритма: при постоянной скорости шума на оракуле квадратичное ускорение Гровера разрушается полностью~\cite{regev2008}. Для умеренных $N$ (обычный поиск) шаг остаётся на порядки выше пола, и механизм неощутим.

%%% ===================================================================
\section{Заключение}
%%% ===================================================================


Привязка флуктуации Каройхази к собственному времени и постулат, что ветви не разделяют общих часов, дают однострочный закон декогеренции (при $\xi=\mathcal O(1)$)
\begin{equation}
\boxed{V/V_0=\exp\!\left[-\left(\frac{mc^2}{\hbar}\right)^{2}\!\bigl(\xi\,t_P^2\,t\bigr)^{2/3}\right],}
\end{equation}
квадратичный по массе, невосприимчивый к электромагнитной экранировке и к неинтерферометрическим границам коллапса; интерферометрия натрия (2026) уже ограничивает $\xi\lesssim0{,}3$, не отвергая модель, но прижимая коэффициент. Сочетание того же $t_{\mathrm{dec}}$ с пределом Марголуса--Левитина даёт максимальное число ортогональных операций $N_{\max}=(2/\pi\xi)(m_P/m)^2$, схлопывающееся до $\mathcal O(1)$ на планковской массе. Его самая смелая черта~--- декогеренция, задаваемая полной энергией ветви, а не разностью энергий,~--- одновременно и его самое уязвимое допущение. Тем не менее модель резко фальсифицируема: единственный ближайший эксперимент на мегадальтонном масштабе либо покажет полосы полного контраста, закрыв предложение, либо покажет их предсказанное гашение с характерным скейлингом $m^2 t^{2/3}$. То, что фундаментальный вопрос о реальности волновой функции сводится к измерению контраста на пучке наночастиц,~--- на наш взгляд, достаточная причина это измерение провести.

\bibliographystyle{unsrt}
\bibliography{refs}

\end{document}
